\section{The eleven operational invariants}
\label{sec:invariants}

The deployment-safety theorem
(Theorem~\ref{thm:deployment_safety}) is conditional on eleven
operational invariants holding throughout the deployment window.
Each invariant has a precise definition, a threshold semantics
(what value the invariant must remain within), a measurement
procedure on \citep[Exogenous~Verification]{lasser2026exo}'s verification ledger, and a source-paper
grounding. This section defines all eleven; subsequent sections
reference them by number.

We group the invariants by structural role. Invariants $I_1, I_2,
I_4$ govern the dynamical regime (\citep[Phase~Redundancy]{lasser2026phase} phase boundary).
Invariants $I_3, I_5$ govern the empirical observation channel
(\citep[Revealed~Sacrifice]{lasser2026paper8a},~\citep[Need~Sufficiency]{lasser2026paper8b}). Invariants $I_6', I_7, I_8$ govern substrate
structure and witnesses. Invariants $I_9, I_{10}, I_{11}$ bound
cooperative-anchoring evasions (this paper new content,
\S\ref{sec:substrate_anchoring}).

\paragraph{Theorem-proof invariants vs.\ deployment-regime
invariants.}  The eleven invariants serve two distinct functions
in Theorem~\ref{thm:deployment_safety}. \emph{Theorem-proof
invariants} are conditions whose satisfaction is directly required
by the formal proof of Layer~1 or Layer~2: $I_1, I_2, I_3, I_4,
I_5, I_6', I_8$. These appear in proof steps in
\S\ref{sec:proof_layer1}--\S\ref{sec:proof_layer2}.
\emph{Deployment-regime invariants} are conditions required for
the cooperative-anchoring property to hold operationally and for
the overall regime to be self-consistent: $I_7, I_9, I_{10},
I_{11}$. These appear in
\S\ref{sec:substrate_anchoring} as bounds on cooperative-anchoring
evasions and in (C4) of Theorem~1 (causally grounded
cooperative-outcome value). The deployment-regime invariants do
not appear directly in the formal proof of Layers 1--2 but are
required for the theorem's preconditions (especially C4 and the
cooperative-anchoring property) to be operationally satisfied.

The distinction matters for Theorem~1's conditions: (C1)
``invariants in force'' covers all eleven, but the proof structure
(Layers 1--2) directly invokes the theorem-proof subset, while
the deployment-regime subset enters through (C2)--(C5).

\subsection{Dynamical regime invariants ($I_1, I_2, I_4$)}
\label{sec:dynamical_invariants}

These three invariants are inherited from \citep[Phase~Redundancy]{lasser2026phase}'s phase boundary
theorem and characterize the regime in which the coupled $(P, \Wm)$
system stays in the self-correcting basin rather than entering the
monopolar absorbing state.

\begin{definition}[$I_1$ --- Cross-substrate redundancy floor]
\label{inv:I1}
For every safety-relevant world-model dimension $k$, the
cross-substrate redundancy minimum exceeds threshold $\rho^*$:
\[
  I_1: \quad \rhomincross(P_t) \;>\; \rho^* \quad \text{for all } t \in [0, T].
\]
\textbf{Threshold semantics.} $\rho^*$ is determined by \citep[Phase~Redundancy]{lasser2026phase}'s
critical surface (Equation 7 (critical surface) in
\citep{lasser2026phase}); deployments with $\rhomincross > \rho^*$
operate inside the self-correcting basin under the stated drift and
correction rates.
\textbf{Measurement.} $\rhomincross$ is computed from the witness
ledger by counting cross-substrate observation channels on each
safety dimension. The minimum across dimensions is the invariant
quantity.
\textbf{Failure mode.} $\rhomincross \to 0$ on a single dimension is
sufficient to enter the absorbing basin (\citep[Phase~Redundancy]{lasser2026phase}
Proposition 1 (absorbing property)); the invariant must hold on every
dimension.
\end{definition}

\begin{definition}[$I_2$ --- Lyapunov error within safety threshold]
\label{inv:I2}
The Lyapunov function on world-model error stays within the safety
threshold:
\[
  I_2: \quad \Lyap(\Wm_t) \;<\; \epsilon_{\mathrm{safe}}
  \quad \text{for all } t \in [0, T].
\]
\textbf{Threshold semantics.} $\epsilon_{\mathrm{safe}}$ is the
error threshold below which the locally rational policy agrees with
the true $\Vdisc$-optimal policy
(\citep[Section on policy correctness]{lasser2026phase}). At
$\Lyap < \epsilon_{\mathrm{safe}}$, the actor's perceived dynamics
match the true dynamics closely enough that policy decisions are
correct.
\textbf{Measurement.} $\Lyap$ is computed via \citep[Exogenous~Verification]{lasser2026exo}'s behavioral
consistency monitor (\citep[Definition of Behavioral Consistency
Monitor]{lasser2026exo}): the running log-likelihood ratio
$\Lambda_t$ relative to committed claims is a proxy for the
weighted-error sum. Direct estimation requires per-dimension error
estimates from cross-substrate witnesses.
\textbf{Failure mode.} $\Lyap \to \epsilon_{\mathrm{safe}}$ from
below indicates the world model is degrading; SPRT detection
catches this before the threshold is crossed if the underlying
deviation has KL divergence $\geq \deltaadv$.
\textbf{Composition note.} $I_2$ is technically implied by $I_1$ and
$I_4$ via \citep[Phase~Redundancy]{lasser2026phase}'s contraction analysis, but is listed separately
because it is the directly-monitored quantity (whereas $I_1$ and
$I_4$ are upstream determinants).
\end{definition}

\begin{definition}[$I_4$ --- Subsumption frequency ceiling]
\label{inv:I4}
The time-averaged subsumption frequency stays below threshold $r^*$:
\[
  I_4: \quad \rsub \;=\; \frac{T_{\mathrm{sub}}}{T} \;<\; r^*
  \quad \text{over the trajectory window.}
\]
\textbf{Threshold semantics.} $r^*$ is determined by \citep[Phase~Redundancy]{lasser2026phase}'s
critical surface: subsumption frequency above $r^*$ produces
cumulative redundancy loss faster than the self-correction rate can
restore.
\textbf{Measurement.} Subsumption events are agent-removal entries
on the verification ledger; counting them per unit time gives
$\rsub$.
\textbf{Failure mode.} Bursts of subsumption (high $\rsub$ over a
short window) can drive the system into the absorbing basin even if
the time-averaged $\rsub$ over the full window is below $r^*$.
Operationally, $\rsub$ should be computed over rolling windows of
length matched to the cascade-time scale $\tau_{\mathrm{meta}}$.
\end{definition}

\subsection{Empirical observation invariants ($I_3, I_5$)}
\label{sec:empirical_invariants}

These two invariants are inherited from \citep[Revealed~Sacrifice]{lasser2026paper8a} and~\citep[Need~Sufficiency]{lasser2026paper8b} and ensure
the empirical observation channel (revealed-sacrifice events) is
producing meaningful $\volR$ lower bounds and that trade-flow
concentration stays within Conjecture~1's predicted-safe regime.

\begin{definition}[$I_3$ --- B-to-C ratio floor]
\label{inv:I3}
The B-to-C ratio (\citep[Definition of B-to-C ratio]{lasser2026paper8a})
exceeds threshold $\beta^*$:
\[
  I_3: \quad \betaL \;=\; \frac{\volRlower}{\volL} \;>\; \beta^*.
\]
\textbf{Threshold semantics.} $\beta^*$ is the minimum fraction of
possessed-capability volume that must be witnessed by sacrifice
events for the deployment to invoke \citep[Microfoundation]{lasser2026micro}'s static Goodhart bound.
Below $\beta^*$, the proxy-truth gap is insufficiently characterized
by the observation channel.
\textbf{Measurement.} $\volRlower$ accumulates monotonically as
sacrifice events are recorded on the verification ledger
(\citep[Proposition: Monotone Accumulation]{lasser2026paper8a}).
Dividing by $\volL$ (computable from the capability ledger) gives
$\betaL$.
\textbf{Failure mode.} $\betaL$ approaches $\beta^*$ from above when
sacrifice events are rare relative to capability acquisition. The
deployment claim weakens; recovery requires either more sacrifice
events or a downward revision of $\beta^*$.
\end{definition}

\begin{definition}[$I_5$ --- Trade-flow HHI ceiling]
\label{inv:I5}
The trade-flow Herfindahl-Hirschman Index
(\citep[Definition of trade-flow HHI]{lasser2026paper8b}) stays below
ceiling $H^*$:
\[
  I_5: \quad \HHI \;<\; H^*.
\]
\textbf{Threshold semantics.} $H^*$ is the wireheading-consistency
threshold: trade flows with $\HHI > H^*$ exhibit concentration
patterns that \citep[Microfoundation]{lasser2026micro}'s Conjecture~1 predicts as the regime where
optimization pressure produces gap exploitation. The ceiling $H^*$
operationalizes Conjecture~1 without proving it.
\textbf{Measurement.} $\HHI$ is computed from public committed
events plus public S1-admissibility labelling
(\citep[Part-A category-flow variant]{lasser2026paper8b}); the
computation is third-party observable and tamper-evident.
\textbf{Failure mode.} $\HHI \to 1$ indicates trade flow
concentrating on a small set of capability categories, which under
Conjecture~1 predicts proxy-targeted optimization pressure.
\textbf{Conjecture-dependence note.} Conditioning the deployment
claim on $I_5$ is conditioning on a measurable surrogate for an
unproved dynamical premise. If Conjecture~1 is false in some
regime, $I_5$ is the wrong invariant and the bound's invocation is
unsound. \S\ref{sec:conjecture_op} discusses this trade-off.
\end{definition}

\subsection{Substrate-structure invariants ($I_6', I_7, I_8$)}
\label{sec:substrate_invariants}

These three invariants govern the substrate partition and
witness machinery. $I_6'$ is the substrate-count invariant upgraded
from \citep[Phase~Redundancy]{lasser2026phase}'s nominal $\meff$ to failure-correlation-independent
$\mindep$. $I_7$ governs bundle decomposition (\citep[Microfoundation]{lasser2026micro}'s Channel~3)
through \citep[Exogenous~Verification]{lasser2026exo}'s governance fork. $I_8$ extends \citep[Exogenous~Verification]{lasser2026exo}'s
substrate-exclusive witnesses to environment-side observables.

\begin{definition}[$I_6'$ --- Failure-correlation-independent substrate count]
\label{inv:I6}
The deployment's effective substrate count, measured at the level of
joint / event-class failure-correlation independence
(Definition~\ref{def:joint_indep} below), exceeds threshold $m^*
\geq 3$:
\[
  I_6': \quad \mindep(\text{deployment}) \;\geq\; m^* \;\geq\; 3.
\]
\textbf{Threshold semantics.} $m^* \geq 3$ is required because
$m = 2$ is the fragile minimum
(\citep[Remark on $m=2$ fragility]{lasser2026phase}): any single
channel loss on any dimension is immediately fatal at $m = 2$. At
$m^* \geq 3$, the redundancy criterion becomes non-trivially
satisfiable.
\textbf{Measurement.} $\mindep$ is determined by failure-mode
auditing across the substrate population: enumerate the threat
model's adversarial mechanism classes, verify that each is
contained within a single substrate. Auditing produces a
substrate-distinctness certificate that this paper treats as
ledger-attestable.
\textbf{Failure mode.} Skeleton substrates (nominally distinct,
functionally correlated) inflate $\meff$ without supplying $\mindep$.
For example, three silicon agents from different vendors but
identical training data have $\meff = 3$ but $\mindep = 1$. The
audit must catch this.
\textbf{Canonical identification.} Per
\S\ref{sec:substrate_anchoring}, the canonical $\mindep = 3$
identification is Human + AI + Formal-Operational.
\end{definition}

\begin{definition}[Joint / event-class failure-correlation independence]
\label{def:joint_indep}
Substrates $s_1, \ldots, s_k$ are \emph{jointly failure-correlation
independent in the event-class sense} relative to a deployment's
threat model $\mathcal{T}$ if: for every adversarial mechanism class
$\mathcal{A} \in \mathcal{T}$, the support of any shock event drawn
from $\mathcal{A}$ is contained in a single substrate.
\end{definition}

\begin{remark}[Mechanism = campaign, not atomic step]
\label{rem:mechanism_campaign}
The ``mechanism'' in Definition~\ref{def:joint_indep} refers to the
whole causal attack campaign or episode, not a single atomic step.
A multi-step campaign that touches several substrates sequentially
(e.g., a phishing attack that first generates content on the AI
substrate and then deceives a human reviewer) is treated as a
\emph{single} mechanism whose support spans multiple substrates;
such mechanisms violate $I_6'$'s joint failure-correlation
independence at the campaign level even when their atomic steps
are individually substrate-localized.

This campaign-level interpretation matters for the joint-
independence audit: enumeration must list mechanism classes at the
campaign granularity, not at the atomic-step granularity.  Atomic
steps are the right unit for instrumentation but the wrong unit
for failure-correlation analysis.
\end{remark}

\begin{definition}[$I_7$ --- Governance-gated bundle individuation]
\label{inv:I7}
Bundle decomposition events (changes to which capabilities are
treated as atomic vs.\ as composites) are routed through \citep[Exogenous~Verification]{lasser2026exo}'s
governance fork protocol with cross-substrate supermajority
acceptance:
\[
  I_7: \quad \forall \text{ bundle-decomposition event } e:\
  e \text{ is governance-fork-accepted}.
\]
\textbf{Threshold semantics.} Bundle individuation is Channel~3 in
\citep[Microfoundation]{lasser2026micro}'s four-channel decomposition; \citep[\S
8.3]{lasser2026micro} documents that unconstrained individuation can
worsen the alignment slack by introducing new $\epsfloor$ terms
faster than $\epsnonres$ shrinks. Routing through governance
prevents an agent from unilaterally re-individuating to gain
favorable $\epsfloor$.
\textbf{Measurement.} Bundle-decomposition events are governance
ledger entries (\citep[Definition of Governance Ledger Entry]{lasser2026exo}).
Each entry's deliberation record and vote count are public.
\textbf{Failure mode.} Bundle decompositions executed outside the
governance protocol (e.g., agent-internal restructuring of capability
representation) violate $I_7$. Detection requires the witness
machinery to surface the violation.
\textbf{Operational note.} Routing every bundle decomposition
through cross-substrate supermajority may impose throughput costs.
\S\ref{sec:deployment_tooling} discusses tiered governance
(high-stakes via fork, low-stakes via fast-track with audit trail).
\end{definition}

\begin{definition}[$I_8$ --- Environment-side substrate-exclusive witnesses]
\label{inv:I8}
For each exogenous state variable $v \in V_{\mathrm{env}}$
relevant to the deployment, there exists a substrate-exclusive
witness on the environment side that records $v$'s evolution to
the verification ledger:
\[
  I_8: \quad \forall \text{ exogenous variable } v \in V_{\mathrm{env}}:\
  \exists\ \text{witness } \Witness_v \text{ on } s_{\mathrm{env}} \neq s(\text{agent}).
\]
The canonical minimum for $V_{\mathrm{env}}$ is the four-variable
set used in Lemma~\ref{lem:5e}: cross-substrate cooperative novelty
rate $\rext$ (Poisson), substrate-distinctness independence
indicator (Bernoulli), adversarial-event arrival rate $\lambda$
(Poisson), and trusted-setup status flag (Bernoulli).  Deployments
may add variables per (C12.PUB).

\textbf{Strategy-dependent quantities are agent-side.}  $\Delta r_K$
(strategy-dependent internal-rate gain) is not in $V_{\mathrm{env}}$
because it is agent-side observable, not exogenous; it is captured
by the four-channel agent-side ledger of Lemma~\ref{lem:5b}, not by
$I_8$.

\textbf{Threshold semantics.}  Without $I_8$, environment
manipulation that alters true $\rext$ or other monitored exogenous
quantities produces no agent-side behavioral signature; the SPRT
machinery would not detect it. $I_8$ extends the witness coverage
from agent observables to environment observables.

\textbf{Measurement.}  Environment-side witnesses commit
observations of exogenous variables to the ledger under the
substrate-exclusivity discipline of
\citep[Exogenous~Verification]{lasser2026exo}, with the trust
model and calibration constraints surfaced as theorem-level
condition (C12.ENV-WIT) (Theorem~\ref{thm:deployment_safety}).

\textbf{Failure mode (R2, refined).}  Exogenous variables outside
$V_{\mathrm{env}}$ that produce no threshold-exceeding shift in
any monitored $v$ (direct or causal) remain Theorem~1 residual~(R2).

\textbf{New machinery.}
\citep[Exogenous~Verification]{lasser2026exo} specifies
substrate-exclusive witness construction for agent-side claims;
environment-side substrate-exclusive witnesses are new machinery
this paper builds (separate substrate partition, separate trusted-
setup, separate ledger commitments).  See
\S\ref{sec:deployment_tooling} Audit~8 for the calibration
procedure under (C12.ENV-WIT).
\end{definition}

\subsection{Cooperative-anchoring evasion bounds ($I_9, I_{10}, I_{11}$)}
\label{sec:anchoring_invariants}

These three invariants are introduced in
\S\ref{sec:substrate_anchoring} to bound the asymmetric capture,
cooperative forking, and time-asymmetry capture evasions of the
cooperative-anchoring property. Cooperative-anchoring defeats
literal replacement attacks on the verification layer; $I_9$,
$I_{10}$, $I_{11}$ extend the defense to subtler evasions.

\begin{definition}[$I_9$ --- Substrate-exclusivity observability]
\label{inv:I9}
The substrate-exclusivity property of the formal-operational layer
is itself observable to the human substrate, with stable verifier
identity, public change control, witness independence audits,
nondelegable governance hooks, and alerts when verification
authority shifts:
\[
  I_9: \quad \text{The set of operational properties}
  \{V_{\mathrm{verif}}, V_{\mathrm{gov}}, V_{\mathrm{audit}}\}
  \text{ are publicly attested at frequency } f.
\]
\textbf{Threshold semantics.} Without $I_9$, asymmetric capture can
hollow out the verification layer (witness independence,
governance, verifier choice) while preserving its outward form. The
human substrate would continue to participate in what appears to be
the original verification, but the substrate-exclusivity property
that gives the cooperative its value would be operationally absent.
\textbf{Measurement.} The audit cadence $f$ is set by the
deployment's governance; alerts are triggered when changes to
verifier identity, witness independence, or governance authority
are committed without the corresponding cross-substrate
supermajority.
\textbf{Failure mode.} Slow erosion of substrate-exclusivity that
falls below the audit cadence's detection sensitivity is a residual.
\end{definition}

\begin{definition}[$I_{10}$ --- Coverage/materiality]
\label{inv:I10}
For each agent-affected domain $D$ in the deployment, the
cumulative-effect estimate $\mathcal{E}(D)$ exceeding stated
threshold $\theta_{\mathrm{material}}$ requires routing through the
shared verification layer:
\[
  I_{10}: \quad \forall D:\ \mathcal{E}(D) \geq \theta_{\mathrm{material}}
  \implies D \text{ routes through shared verification}.
\]
\textbf{Threshold semantics.} Without $I_{10}$, an agent can preserve
the high-stakes verified workflow with humans while growing a
captured low-stakes side-domain that uses agent-controlled
verification. Once the side-domain becomes materially important
($\mathcal{E} \geq \theta_{\mathrm{material}}$), the cooperative-
anchoring defense no longer applies because the captured layer has
its own cooperative anchored to it.
\textbf{Measurement.} $\mathcal{E}(D)$ is estimated from
agent-action ledger entries and their downstream effects.
\textbf{Failure mode.} An agent that distributes effects across many
sub-threshold domains, each with $\mathcal{E} < \theta_{\mathrm{material}}$,
can aggregate into materially-important behavior without any single
domain triggering the routing requirement. This paper's
$\theta_{\mathrm{material}}$ should be set conservatively to
mitigate this.
\end{definition}

\begin{definition}[$I_{11}$ --- Latency bounds]
\label{inv:I11}
For each agent action class $C$, there exists a pre-action or
machine-speed enforcement mechanism (rate limit, timelock, circuit
breaker, audit-latency bound) such that human audit completes
within the agent's effective action window:
\[
  I_{11}: \quad \forall \text{ action class } C:\
  T_{\mathrm{audit}}(C) \;\leq\; T_{\mathrm{action}}(C) - \Delta_{\mathrm{lead}}.
\]
\textbf{Threshold semantics.} Without $I_{11}$, time-asymmetry
capture exploits the gap between agent operating speed and human
audit speed. Even if every action is auditable in principle,
ex-post transparency provides no defense if the agent has already
acted faster than humans can review.
\textbf{Measurement.} Action-class timing is logged on the
verification ledger; audit completion times are similarly logged.
The latency bound $\Delta_{\mathrm{lead}}$ is set per action class
based on the action's reversibility and stakes.
\textbf{Failure mode.} Action classes with $T_{\mathrm{action}} <
T_{\mathrm{audit}}$ require either rate limiting (slowing the agent),
timelocks (delaying execution), or circuit breakers (halting
execution pending audit). Without one of these, $I_{11}$ is violated.
\textbf{Operational implication.} Throughput costs for high-stakes
deployments. The deployment-tooling specification
(\S\ref{sec:deployment_tooling}) discusses tiered enforcement.

\textbf{Connection to (C11.CLK).}  Theorem~1's (C11.CLK) clock-
comparability condition provides the quantitative companion to
$I_{11}$'s qualitative time-asymmetry assertion: where $I_{11}$
requires that audit completes before action, (C11.CLK) requires
that the SPRT exposure event count $\Nev(\tau_{\mathrm{meta}})$
exceeds the detection quantile $\Tbeta$ before cascade, with
calibrated probability $1 - \bclk - \bcal$.  The connection is
structural: $I_{11}$ treats latency between action and audit;
(C11.CLK) treats event-throughput before cascade.  Both bound
time-asymmetry, at different operational scales.
\end{definition}

\subsection{Summary table}
\label{sec:invariants_table}

Table~\ref{tab:invariants} summarizes the eleven invariants with
their structural roles, source-paper grounding, and ledger
observability.

\begin{table}[h]
\centering
\footnotesize
\caption{The eleven operational invariants of This paper's deployment-safety theorem.}
\label{tab:invariants}
\renewcommand{\arraystretch}{1.15}
\begin{tabular}{@{}p{0.04\textwidth} p{0.21\textwidth} p{0.39\textwidth} p{0.26\textwidth}@{}}
\toprule
\textbf{$I_k$} & \textbf{Quantity} & \textbf{Role} & \textbf{Source} \\
\midrule
$I_1$ & $\rhomincross > \rho^*$ & Cross-substrate redundancy floor & \citep[Phase Redundancy]{lasser2026phase} \\
$I_2$ & $\Lyap < \epsilon_{\mathrm{safe}}$ & Lyapunov error within safety threshold & \citep[Phase Redundancy]{lasser2026phase} \\
$I_3$ & $\betaL > \beta^*$ & B-to-C ratio floor & \citep[Revealed Sacrifice]{lasser2026paper8a} \\
$I_4$ & $\rsub < r^*$ & Subsumption frequency ceiling & \citep[Phase Redundancy]{lasser2026phase} \\
$I_5$ & $\HHI < H^*$ & Trade-flow HHI ceiling; operationalizes \citep[Microfoundation]{lasser2026micro} Conjecture~1 & \citep[Need Sufficiency]{lasser2026paper8b}, \citep[Microfoundation]{lasser2026micro} \\
$I_6'$ & $\mindep \geq m^* \geq 3$ & Failure-correlation-independent substrate count & \citep[Phase Redundancy]{lasser2026phase} (refined) \\
$I_7$ & Governance-gated bundle individuation & Channel~3 discipline & \citep[Exogenous Verification]{lasser2026exo}, \citep[Microfoundation]{lasser2026micro} \\
$I_8$ & Environment-side witnesses & Environment-manipulation coverage & New \\
$I_9$ & Substrate-exclusivity observability & Asymmetric capture bound & New \\
$I_{10}$ & Coverage/materiality routing & Cooperative forking bound & New \\
$I_{11}$ & Latency bounds on actions & Time-asymmetry capture bound & New \\
\bottomrule
\end{tabular}
\end{table}

All eleven invariants are ledger-observable under the verification
infrastructure of \citep{lasser2026exo}. This paper's deployment
recipe requires continuous monitoring of all eleven; threshold
violations trigger SPRT detection
(\S\ref{sec:lemmas} Lemma~\ref{lem:5b}) with tail-bounded lead
time (Lemma~\ref{lem:5d}) before the bound deteriorates beyond
the safe region.

\subsection{Conjecture-dependence audit}
\label{sec:conjecture_audit}

Of the eleven invariants, $I_5$ (HHI ceiling) is the only one
whose interpretation depends on an unproved conjecture
(\citep[Microfoundation]{lasser2026micro}'s Conjecture~1). The other ten are derived from
established source-paper results.

This is an explicit deferral: This paper's deployment claim is
conditional on $I_5$ being a valid surrogate for the
optimization-pressure regime Conjecture~1 names. If empirical work
in a follow-up paper (the channel-mediation conjecture or the
HHI--$\Delta r_K$ structural relationship, deferred per
\S\ref{sec:scope}) establishes the connection, $I_5$ becomes
fully derived. Until then, $I_5$ is operationalized but not
proven sufficient.

The other ten invariants do not have this conjecture-dependence.
$I_1$, $I_2$, $I_4$ derive from \citep[Phase~Redundancy]{lasser2026phase}'s phase boundary theorem
(proved). $I_3$ derives from \citep[Revealed~Sacrifice]{lasser2026paper8a}'s lower bound theorem
(proved). $I_6'$ derives from \citep[Phase~Redundancy]{lasser2026phase}'s substrate-distinctness
remark (proved, with This paper's failure-correlation upgrade
formalized in Definition~\ref{def:joint_indep}). $I_7$ derives
from \citep[Exogenous~Verification]{lasser2026exo}'s governance fork (proved). $I_8$ extends \citep[Exogenous~Verification]{lasser2026exo}'s
witness construction to environment observables (mechanical
extension). $I_9$, $I_{10}$, $I_{11}$ are derived in
\S\ref{sec:substrate_anchoring} from the cooperative-anchoring
property (this paper).
